\chapter{数学分析}
\label{chap:mathematical-analysis}

第~\ref{chap:sets-functions-and-proofs}章说明了怎样写清对象、条件与保证。本章进一步追问： 当计算过程包含无限多步，或者函数需要在连续区域上取平均时，有限步骤中的代数规则还能否直接沿用？答案取决于极限。导数用极限描述函数在一点附近的变化，积分用极限汇总一个区域上的贡献，迭代算法也常以极限作为最终结果。极限存在并不意味着它能与求导、积分或另一重极限交换；每一种交换都需要相应的控制条件。

本章假定读者已经学过一元微积分和基本矩阵运算，从这些熟悉的计算出发补齐后续机器学习推导所需的条件。我们先研究数列与连续性，明确无限过程何时落到允许的空间中；随后把导数组织成局部线性映射，并建立向量和矩阵微分；再由有限加权和进入一般积分，分析换元、积分次序和参数求导。函数不光滑或函数列发生变化时，还必须说明控制的是函数值、导数、平均误差还是最坏输入误差。章末的压缩映射定理把这些线索收拢起来：它不仅保证迭代极限存在， 还给出唯一性和可核对的误差界。
\section{数列、级数与误差趋零}
\label{sec:analysis-sequences-series-errors}
\subsection{极限描述最终稳定到哪里}

\term{数列}（sequence）是以自然数为索引的一族数，写作$(a_n)_{n\geq 1}$。若对每个$\varepsilon>0$，都存在只依赖于$\varepsilon$的整数$N$，使得
\[   n\geq N\quad\Longrightarrow\quad |a_n-a|<\varepsilon, \]
就称$a_n$\term{收敛}（converge）到$a$，记作$a_n\to a$。这一定义要求同一个 $N$控制其后的全部项，而不是只要求数列偶尔靠近$a$。极限若存在便唯一：若$a\neq b$ 都是极限，取$\varepsilon=|a-b|/3$，充分靠后的$a_n$会同时落入两个不相交的邻域， 产生矛盾。

收敛数列必有界。事实上，取$\varepsilon=1$后，除有限个前项外都有 $|a_n|\leq |a|+1$，再把有限个前项的绝对值一并取最大值即可。反向结论通常不成立， 例如$a_n=(-1)^n$有界却不收敛。若再加入单调性，情况便不同。

这里需要使用实数的\term{完备性}（completeness）。它的一种等价表述是确界公理：
每个非空且有上界的实数集合都有实数上确界。稍后讨论柯西数列时，还会看到完备性的
另一种等价表述。
\begin{theorem}[单调有界收敛定理]
\label{thm:analysis-monotone-bounded-sequence}
若实数列$(a_n)$单调不减且有上界，则$a_n$收敛到 $a=\sup\{a_n:n\in\mathbb{N}\}$。单调不增且有下界的情形类似。
\end{theorem}
\begin{proof}
这里只证明单调不减的情形。集合$\{a_n:n\in\mathbb{N}\}$非空且有上界，实数的完备性保证其上确界$a$存在。给定$\varepsilon>0$，数$a-\varepsilon$不可能仍是上界，所以存在$N$使$a_N>a-\varepsilon$。当$n\geq N$时，单调性给出
\[   a-\varepsilon<a_N\leq a_n\leq a, \]
从而$|a_n-a|<\varepsilon$。这正是$a_n\to a$。
\end{proof}

定理中的“有界”与“单调”缺一不可：$n$单调却无上界，$(-1)^n$有界却不单调。证明还使用了实数的完备性；若只在有理数中考察逐步逼近$\sqrt{2}$的有理数列，所有项都在有理数中，极限却不在这个空间中。第~\ref{sec:analysis-complete-contraction}节会把这种“柯西数列的极限仍留在空间内”的性质推广到一般度量空间。
\subsection{柯西条件只比较数列内部}

有时我们还不知道候选极限，却能证明后面的项彼此越来越近。\term{柯西数列}
（Cauchy sequence）满足
\[   \forall\varepsilon>0,\quad   \exists N,\quad   \forall m,n\geq N,\quad |a_m-a_n|<\varepsilon. \]
每个收敛实数列都是柯西数列，因为 $|a_m-a_n|\leq|a_m-a|+|a_n-a|$。反过来，每个实柯西数列也收敛；这是实数完备性的等价表述之一。柯西条件的价值在于，它只使用已经算出的项，适合分析尚不知道极限为何物的迭代。

若能控制相邻差，还可以借助三角不等式验证柯西条件。对$m>n$，
\begin{equation}
  |a_m-a_n|
  \leq
  \sum_{k=n}^{m-1}|a_{k+1}-a_k|.
\label{eq:analysis-cauchy-tail-bound}
\end{equation}
因此，只要相邻差的无穷级数收敛，其尾和就趋于零，$(a_n)$必为柯西数列。章末证明压缩迭代收敛时，关键步骤正是把相邻迭代差控制成可求和的几何序列。
\subsection{级数控制累计量}

\term{无穷级数}（infinite series）$\sum_{n=1}^{\infty}a_n$是否收敛，是指部分和 $s_N=\sum_{n=1}^{N}a_n$是否有有限极限。级数收敛必然推出$a_n=s_n-s_{n-1}\to0$， 但单项趋零并不足够。例如调和级数$\sum_{n=1}^{\infty}1/n$发散。把从$2^k$到 $2^{k+1}-1$的项分成一组，该组有$2^k$项，每项至少为$1/2^{k+1}$，所以每组之和至少为$1/2$；部分和因而可以超过任意给定常数。

当$|r|<1$时，有限几何和
\[   \sum_{n=0}^{N}r^n=\frac{1-r^{N+1}}{1-r} \]
在$N\to\infty$时给出
\begin{equation}
  \sum_{n=0}^{\infty}r^n=\frac{1}{1-r}.
\label{eq:analysis-geometric-series}
\end{equation}
$|r|<1$不可删去：$r=1$时部分和不断增长，$r=-1$时部分和振荡，$|r|>1$时单项甚至不趋于零。折扣回报和几何误差界都依赖这一条件。

若$\sum_n|a_n|<\infty$，就称级数\term{绝对收敛}（absolutely converge）。由式~\eqref{eq:analysis-cauchy-tail-bound}对应的尾和估计，绝对收敛蕴含收敛。绝对收敛还允许较自由地重新排列求和次序；只有条件收敛时，重排可能改变和甚至使其发散。第~\ref{sec:analysis-limit-integral-interchange}节中的Fubini定理，是这一原则在积分上的对应版本。
\subsection{单步、累计与可求和误差}

设第$t$步计算产生非负误差$\delta_t$。“单步误差趋零”只表示$\delta_t\to0$； “前$T$步累计误差”是$\sum_{t=0}^{T-1}\delta_t$；“误差可求和”则要求 $\sum_{t=0}^{\infty}\delta_t<\infty$。三者不能混用。例如 $\delta_t=1/(t+1)$逐步趋零，但累计误差无界；$\delta_t=1/(t+1)^2$才可求和。若误差进入第~\ref{chap:sets-functions-and-proofs}章式
\eqref{eq:prelim-perturbed-recurrence}那样的稳定递推，较早误差还会被几何权重衰减， 此时应分析加权累计量，而不是只看$\delta_t$本身。

步长序列也常同时承担“持续前进”和“抑制平方误差”两项职责。典型条件
\[   \sum_{t=1}^{\infty}\alpha_t=\infty,   \qquad   \sum_{t=1}^{\infty}\alpha_t^2<\infty \]
可由$\alpha_t=1/t$满足：第一项是调和级数，第二项是收敛的$p$级数。这里尚不能据此推出随机算法收敛，因为还缺少噪声和目标函数的条件；这个例子只说明两个累计对象可以有不同结论。后续分析折扣、步长和扰动时，都应先问清趋零的是哪一项、累计采用什么权重， 以及所需级数是否真的收敛。
\section{距离、连续性与极值存在}
\label{sec:analysis-metric-continuity-extrema}
\subsection{距离把“靠近”推广到一般对象}

实数上的绝对值能够衡量两点距离，但函数、概率分布或离散状态未必适合用绝对值比较。集合$X$上的\term{度量}（metric）是函数$d:X\times X\to[0,\infty)$，满足对任意 $x,y,z\in X$，
\[   d(x,y)=0\Longleftrightarrow x=y,\qquad   d(x,y)=d(y,x),\qquad   d(x,z)\leq d(x,y)+d(y,z). \]
带有度量的集合$(X,d)$称为\term{度量空间}（metric space）。欧氏空间使用 $d(\bm{x},\bm{y})=\lVert\bm{x}-\bm{y}\rVert_2$；有限集合可以使用离散度量 $d(x,y)=\mathbb{I}\{x\neq y\}$。同一个集合选择不同度量，会改变何谓收敛以及函数是否连续，所以讨论这些性质时必须说明所用距离。

以$x\in X$为中心、半径$r>0$的开球记为
$B(x,r)=\{y\in X:d(x,y)<r\}$。数列$(x_n)$在$(X,d)$中收敛到$x$，是指
\[
  \forall\varepsilon>0,\quad
  \exists N,\quad
  n\geq N\Longrightarrow d(x_n,x)<\varepsilon;
\]
等价地，$d(x_n,x)\to0$。数列$(x_n)$称为柯西数列，是指
\[
  \forall\varepsilon>0,\quad
  \exists N,\quad
  m,n\geq N\Longrightarrow d(x_m,x_n)<\varepsilon.
\]
这两条定义分别把实数中的$|a_n-a|$与$|a_m-a_n|$替换为一般距离。

集合$U\subseteq X$若对每个$x\in U$都包含某个开球$B(x,r)$，就称为
\term{开集}（open set）；集合$C$若其补集$X\setminus C$为开集，就称为
\term{闭集}（closed set）。闭集的一个实用判据是：若$C$中的任意收敛序列
$x_n\to x$都有$x\in C$，则$C$闭。闭与有界描述不同性质；例如$(0,1)$有界但不闭，
$\mathbb{R}$闭但不有界。
\subsection{连续性保证小输入扰动只产生小输出扰动}

设$(X,d_X)$与$(Y,d_Y)$为度量空间。函数$f:X\to Y$在$x\in X$处
\term{连续}（continuous），是指
\[   \forall\varepsilon>0,\quad   \exists\delta>0,\quad   d_X(x,x')<\delta   \Longrightarrow   d_Y\bigl(f(x),f(x')\bigr)<\varepsilon. \]
这里$\delta$可以依赖$x$和$\varepsilon$。等价地，只要$x_n\to x$，就有 $f(x_n)\to f(x)$。连续函数因而允许把一个已经存在的极限送入函数： $f(\lim_n x_n)=\lim_n f(x_n)$。

若对每个$\varepsilon>0$都能找到一个与具体位置无关的$\delta>0$，使任意
$x,x'\in X$只要$d_X(x,x')<\delta$便有
$d_Y(f(x),f(x'))<\varepsilon$，就称$f$\term{一致连续}
（uniformly continuous）。连续性允许不同位置使用不同的$\delta$，一致连续则要求同一个
$\delta$同时控制整个定义域。连续函数在紧致定义域上一致连续，但定义域无界时未必如此：
$f(x)=x^2$在$\mathbb{R}$上连续，取$x'=x+1/x$可见输入距离趋于零而输出距离
$|f(x')-f(x)|=2+1/x^2$不趋于零。

若存在常数$L\geq0$，使所有$x,x'\in X$都满足
\begin{equation}
  d_Y\bigl(f(x),f(x')\bigr)\leq Ld_X(x,x'),
\label{eq:analysis-lipschitz-condition}
\end{equation}
则称$f$为$L$-\term{Lipschitz连续}（Lipschitz continuous）。它不仅用同一个
$\delta$控制整个定义域，还给出输出变化相对于输入变化的线性上界，因而蕴含一致连续。
函数$f(x)=\sqrt{x}$在$[0,1]$上一致连续，却不是Lipschitz连续：若存在有限$L$，
令$x'=0$便会要求$\sqrt{x}\leq Lx$，即$1/\sqrt{x}\leq L$，这在$x\downarrow0$时不可能成立。
\subsection{紧致性把局部控制变成整体控制}

设$K\subseteq X$。若一族开集的并包含$K$，就称这族开集为$K$的
\term{开覆盖}（open cover）；若从中选出的有限个开集仍覆盖$K$，它们就构成
\term{有限子覆盖}（finite subcover）。集合$K$的每个开覆盖都能选出有限子覆盖时，
称$K$为\term{紧致}（compact）。在度量空间中，这一定义等价于
\term{序列紧致性}：集合中的每个序列都有一个子序列收敛到仍在该集合中的点。
后一种表述更适合分析算法产生的点列，因为它直接提供候选极限。

在有限维欧氏空间$\mathbb{R}^d$中，Heine--Borel定理说明紧致等价于闭且有界。这个结论不能原样搬到一般度量空间。考虑所有有界实数序列组成的空间
\[   \ell^\infty   =   \left\{\bm{x}=(x_1,x_2,\ldots):   \sup_{k\geq1}|x_k|<\infty\right\},   \qquad   d(\bm{x},\bm{y})=\sup_{k\geq1}|x_k-y_k|. \]
令$\bm{e}_n$只有第$n$个坐标为$1$，其余坐标为$0$。集合 $\{\bm{e}_n:n\geq1\}$有界，而且在$\ell^\infty$中是闭集；但$n\neq m$时 $d(\bm{e}_n,\bm{e}_m)=1$，所以任何子序列都不是柯西数列，更不可能收敛。一般空间中 “闭且有界”不足以替代紧致，正是因为无穷多个彼此分离的方向可能留在同一个有界球内。
\begin{theorem}[紧致集上的极值定理]
\label{thm:analysis-extreme-value}
设$K$是非空紧致度量空间，$f:K\to\mathbb{R}$连续，则$f$在$K$上同时达到最小值与最大值。
\end{theorem}
\begin{proof}
只证明最小值情形。令$m=\inf_{x\in K}f(x)$。先说明$m$有限：若$f$无下界，可为每个 $n$选取$x_n\in K$使$f(x_n)<-n$。紧致性给出收敛子序列 $x_{n_k}\to x^\star\in K$，连续性却要求$f(x_{n_k})\to f(x^\star)\in\mathbb{R}$， 与$f(x_{n_k})<-n_k\to-\infty$矛盾。

按下确界定义，对每个$n$可选$x_n\in K$使$f(x_n)<m+1/n$。再次由紧致性，存在子序列$x_{n_k}\to x^\star\in K$。连续性给出
\[   f(x^\star)   =\lim_{k\to\infty}f(x_{n_k})   =m, \]
故$x^\star$达到最小值。对$-f$应用同一结论便得到最大值。
\end{proof}

这个证明展示了紧致性的两项作用：从近似最优点列中抽出收敛子序列，并保证极限仍在可行集合内。连续性随后把函数值极限传到该点。任一条件缺失都可能使最优解不存在。
\begin{example}[极值存在条件怎样失效]
\label{ex:analysis-extrema-failures}
在开区间$(0,1)$上，连续函数$f(x)=x$的下确界为$0$，却没有最小点；定义域不紧致， 极限点$0$被排除。把定义域改成闭但不紧致的$\mathbb{R}$，连续函数$f(x)=e^x$仍有下确界$0$而不达到它。另一方面，在紧致集合$[0,1]$上定义
\[   g(0)=1,\qquad g(x)=x\quad(x>0), \]
则$\inf g=0$也未达到；这次失败的是$g$在$0$处不连续。三个例子分别表明“连续”、 “闭”或“有界”中的单独一项都不是一般的极值存在保证。
\end{example}
\section{导数与线性近似}
\label{sec:analysis-derivatives-linear-approximation}
\subsection{导数是局部变化率}

一元函数$f:\mathbb{R}\to\mathbb{R}$在$x$处的\term{导数}（derivative）定义为
\[   f'(x)=\lim_{h\to0}\frac{f(x+h)-f(x)}{h}, \]
前提是这个双侧极限存在。它说明输入发生小变化$h$时，函数增量的一阶部分是$f'(x)h$：
\begin{equation}
  f(x+h)=f(x)+f'(x)h+o(|h|).
\label{eq:analysis-scalar-first-order}
\end{equation}
式中的小$o$比“$h$很小”更精确：余项$r(h)$满足$r(h)/|h|\to0$。导数存在必然推出函数在该点连续，但连续不保证可导，例如$|x|$在$0$处连续，左右导数却分别为 $-1$和$1$。
跳变还会破坏连续性，因而更不可能可导：阈值函数
$\mathbb{I}\{x\geq0\}$在$0$处的左右函数值极限不同。尖点与跳变都会使经典导数失效，
但前者仍连续，后者连零阶近似也无法通过同一个函数值连接两侧。

对$f:\mathbb{R}^d\to\mathbb{R}$，沿方向$\bm{v}\in\mathbb{R}^d$的
\term{方向导数}（directional derivative）是
\[   D_{\bm{v}}f(\bm{x})   =   \lim_{t\to0}   \frac{f(\bm{x}+t\bm{v})-f(\bm{x})}{t}. \]
取$\bm{v}$为第$j$个坐标方向$\bm{e}_j$，就得到\term{偏导数}
（partial derivative）$\partial f/\partial x_j$。方向导数只检查一条直线，而整体可微性要求所有小扰动共享同一个线性近似。
\subsection{全微分统一所有方向}

函数$f:\mathbb{R}^d\to\mathbb{R}$在$\bm{x}$处\term{可微}
（differentiable），是指存在一个线性映射 $A:\mathbb{R}^d\to\mathbb{R}$，使
\begin{equation}
  f(\bm{x}+\bm{h})
  =
  f(\bm{x})+A(\bm{h})+o(\lVert\bm{h}\rVert_2).
\label{eq:analysis-total-differentiability}
\end{equation}
在欧氏空间中，每个这样的线性泛函都可写为 $A(\bm{h})=\langle\nabla f(\bm{x}),\bm{h}\rangle$，其中
\[   \nabla f(\bm{x})   =   \left(     \frac{\partial f}{\partial x_1}(\bm{x}),\ldots,     \frac{\partial f}{\partial x_d}(\bm{x})   \right)^{\top} \]
称为\term{梯度}（gradient）。因此梯度不是若干偏导数的随意排列，而是表示局部线性映射的向量；扰动$\bm{h}$通过内积变成函数的一阶变化。

可微必然推出所有方向导数存在，并且 $D_{\bm{v}}f(\bm{x})=\langle\nabla f(\bm{x}),\bm{v}\rangle$。反向结论不成立， 仅有偏导数更不够。
\begin{example}[偏导数存在但不可微]
\label{ex:analysis-partials-not-differentiable}
定义$f:\mathbb{R}^2\to\mathbb{R}$为
\[   f(x,y)=
\begin{cases}     \dfrac{xy}{\sqrt{x^2+y^2}}, & (x,y)\neq(0,0),\\[6pt]     0, & (x,y)=(0,0).
\end{cases} \]
沿两条坐标轴，函数恒为$0$，所以 $\partial f/\partial x(0,0)=\partial f/\partial y(0,0)=0$。若$f$在原点可微， 候选线性部分只能是零映射。可是沿$\bm{h}=(t,t)$，
\[   \frac{|f(t,t)-f(0,0)|}{\lVert(t,t)\rVert_2}   =   \frac{|t|/\sqrt{2}}{\sqrt{2}|t|}   =\frac12, \]
余项与$\lVert\bm{h}\rVert_2$之比不趋于零，故$f$在原点不可微。偏导数只看坐标轴， 遗漏了其他接近路径。
\end{example}

一个常用的充分条件是：若$f$的所有一阶偏导数在$\bm{x}$的某个邻域内存在，并在 $\bm{x}$处连续，则$f$在$\bm{x}$处可微。注意这只是方便核对的充分条件，不是必要条件； 可微函数的偏导数不一定在该点附近连续。
\subsection{中值定理把局部导数连接到有限变化}

若$f:[a,b]\to\mathbb{R}$在$[a,b]$上连续、在$(a,b)$上可导，则
\term{中值定理}（mean value theorem）保证存在$c\in(a,b)$使
\begin{equation}
  f(b)-f(a)=f'(c)(b-a).
\label{eq:analysis-mean-value-theorem}
\end{equation}
所以，若区间上$|f'(x)|\leq L$，便有$|f(b)-f(a)|\leq L|b-a|$。这说明有界导数可以推出Lipschitz控制。多元标量函数也可沿线段使用一元中值定理：当连接 $\bm{x}$与$\bm{y}$的线段位于定义域且$f$在线段邻域可微时，
\[   |f(\bm{y})-f(\bm{x})|   \leq   \sup_{0\leq t\leq1}   \lVert\nabla f(\bm{x}+t(\bm{y}-\bm{x}))\rVert_2   \lVert\bm{y}-\bm{x}\rVert_2. \]
这里控制的是标量函数值差。对向量值函数，不能一般地断言存在一个共同的$c$让整个向量增量等于单个Jacobian矩阵乘以输入增量；应使用积分形式或范数上界。
\section{向量值函数与链式法则}
\label{sec:analysis-vector-chain-matrix-calculus}
\subsection{Jacobian矩阵表示局部线性映射}

设$f:\mathbb{R}^d\to\mathbb{R}^m$。若存在矩阵 $\bm{J}_f(\bm{x})\in\mathbb{R}^{m\times d}$使
\begin{equation}
  f(\bm{x}+\bm{h})
  =
  f(\bm{x})+\bm{J}_f(\bm{x})\bm{h}
  +o(\lVert\bm{h}\rVert_2),
\label{eq:analysis-jacobian-linearization}
\end{equation}
则称$f$在$\bm{x}$处可微，$\bm{J}_f(\bm{x})$称为
\term{Jacobian矩阵}（Jacobian matrix）。其第$i,j$个元素是 $\partial f_i/\partial x_j$。矩阵有$m$行、$d$列，因为它把$d$维输入扰动映射成 $m$维输出扰动。

若$g:\mathbb{R}^m\to\mathbb{R}^p$在$f(\bm{x})$处可微，则复合映射满足
\begin{equation}
  \bm{J}_{g\circ f}(\bm{x})
  =
  \bm{J}_g\bigl(f(\bm{x})\bigr)\bm{J}_f(\bm{x}).
\label{eq:analysis-vector-chain-rule}
\end{equation}
这就是\term{链式法则}（chain rule）。维度同时给出乘法次序： $(p\times m)(m\times d)=p\times d$，反向相乘往往没有定义。即使两个方阵能够交换书写顺序，矩阵乘法通常也不交换，不能据标量求导经验调换因子。

当最终输出$\ell:\mathbb{R}^m\to\mathbb{R}$为标量时，采用列向量梯度可把链式法则写成
\[   \nabla_{\bm{x}}(\ell\circ f)(\bm{x})   =   \bm{J}_f(\bm{x})^{\top}   \nabla_{\bm{y}}\ell(\bm{y})\big|_{\bm{y}=f(\bm{x})}. \]
局部线性映射在前向计算中按$\bm{J}_f$作用，在标量目标的反向计算中按转置 $\bm{J}_f^{\top}$传回。这是后续反向传播的线性代数核心，但具体网络结构留到模型章节。

例如，若$f:\mathbb{R}^2\to\mathbb{R}^3$且
$g:\mathbb{R}^3\to\mathbb{R}$，则
$\bm{J}_f\in\mathbb{R}^{3\times2}$、$\bm{J}_g\in\mathbb{R}^{1\times3}$，
复合映射的Jacobian形状为
$(1\times3)(3\times2)=1\times2$。对应的列向量梯度从标量输出传回时，
先乘$\bm{J}_g^{\top}$得到三维向量，再乘$\bm{J}_f^{\top}$得到二维向量。
这两个方向描述的是同一个线性复合，只是向量表示的方向相反。

另一个常用例子是$n$维\term{归一化指数映射}（softmax）。对
$\bm{z}\in\mathbb{R}^n$，令
\[
  s_i(\bm{z})
  =
  \frac{e^{z_i}}{\sum_{k=1}^n e^{z_k}},
  \qquad i=1,\ldots,n.
\]
记分母为$Z=\sum_k e^{z_k}$。指数函数求导和商法则给出
\[
  \frac{\partial s_i}{\partial z_j}
  =
  \frac{\mathbb{I}\{i=j\}e^{z_i}Z-e^{z_i}e^{z_j}}{Z^2}
  =
  s_i\bigl(\mathbb{I}\{i=j\}-s_j\bigr).
\]
因此整个Jacobian矩阵为
\begin{equation}
  \bm{J}_{\bm{s}}(\bm{z})
  =
  \operatorname{diag}(\bm{s})-\bm{s}\bm{s}^{\top}.
\label{eq:analysis-softmax-jacobian}
\end{equation}
对角项$s_i(1-s_i)$描述提高$z_i$对自身分量的影响，非对角项$-s_is_j$
来自共同分母：提高一个坐标会压低其余归一化分量。每行元素之和为零，
这也与给所有$z_i$同时加同一常数不改变输出相符。
\subsection{全微分处理矩阵变量}

矩阵$\bm{X}\in\mathbb{R}^{m\times n}$可以看成$mn$个坐标组成的对象。对标量函数 $f(\bm{X})$，本书使用Frobenius迹内积
\[   \langle\bm{A},\bm{B}\rangle_F   =   \operatorname{tr}(\bm{A}^{\top}\bm{B}) \]
定义矩阵梯度：若一阶全微分能够写成
\begin{equation}
  \mathrm{d}f
  =
  \langle\nabla_{\bm{X}}f,\mathrm{d}\bm{X}\rangle_F
  =
  \operatorname{tr}\bigl((\nabla_{\bm{X}}f)^{\top}
  \mathrm{d}\bm{X}\bigr),
\label{eq:analysis-matrix-gradient-definition}
\end{equation}
则$\nabla_{\bm{X}}f$与$\bm{X}$形状相同。计算时可先用乘积法则展开微分，再利用
\[   \operatorname{tr}(\bm{A}\bm{B})   =\operatorname{tr}(\bm{B}\bm{A}),\qquad   \operatorname{tr}(\bm{A}\bm{B}\bm{C})   =\operatorname{tr}(\bm{B}\bm{C}\bm{A}), \]
把$\mathrm{d}\bm{X}$移动到末端。迹只允许循环移位，不能任意交换相邻矩阵。
\subsection{四个常用矩阵微分}

先考虑二次型$q(\bm{x})=\bm{x}^{\top}\bm{A}\bm{x}$，其中 $\bm{A}\in\mathbb{R}^{d\times d}$固定。由乘积法则，
\begin{align*}
  \mathrm{d}q
  &=
  (\mathrm{d}\bm{x})^{\top}\bm{A}\bm{x}
  +\bm{x}^{\top}\bm{A}\,\mathrm{d}\bm{x}\\
  &=
  \bm{x}^{\top}(\bm{A}^{\top}+\bm{A})\,\mathrm{d}\bm{x}.
\end{align*}
因此
\begin{equation}
  \nabla_{\bm{x}}
  \bigl(\bm{x}^{\top}\bm{A}\bm{x}\bigr)
  =
  (\bm{A}+\bm{A}^{\top})\bm{x}.
\label{eq:analysis-quadratic-form-gradient}
\end{equation}
只有当$\bm{A}$对称时才能化为$2\bm{A}\bm{x}$。直接写成$2\bm{A}\bm{x}$会悄然加入对称性条件。

对固定的$\bm{A}\in\mathbb{R}^{m\times n}$，令 $f(\bm{X})=\operatorname{tr}(\bm{A}^{\top}\bm{X})$。它的微分已经符合式
\eqref{eq:analysis-matrix-gradient-definition}，所以
\begin{equation}
  \nabla_{\bm{X}}\operatorname{tr}(\bm{A}^{\top}\bm{X})
  =\bm{A}.
\label{eq:analysis-trace-linear-gradient}
\end{equation}
这个简单关系常用于从迹表达式中读出梯度。

若方阵$\bm{X}$可逆，由恒等式$\bm{X}\bm{X}^{-1}=\bm{I}$取微分可得
\[   (\mathrm{d}\bm{X})\bm{X}^{-1}   +\bm{X}\,\mathrm{d}(\bm{X}^{-1})=\bm{0}. \]
左乘$\bm{X}^{-1}$后得到
\begin{equation}
  \mathrm{d}(\bm{X}^{-1})
  =
  -\bm{X}^{-1}(\mathrm{d}\bm{X})\bm{X}^{-1}.
\label{eq:analysis-inverse-differential}
\end{equation}
两个$\bm{X}^{-1}$位于扰动两侧，不能合并成 $-\bm{X}^{-2}\mathrm{d}\bm{X}$，因为$\mathrm{d}\bm{X}$未必与$\bm{X}^{-1}$交换。

最后设$\bm{X}\in\mathbb{R}^{d\times d}$正定，因而$\det\bm{X}>0$且 $\log\det\bm{X}$有定义。Jacobi行列式公式给出
\[   \mathrm{d}\det\bm{X}   =   \det(\bm{X})\operatorname{tr}   (\bm{X}^{-1}\mathrm{d}\bm{X}). \]
再应用标量对数的链式法则，
\begin{equation}
  \mathrm{d}\log\det\bm{X}
  =
  \operatorname{tr}(\bm{X}^{-1}\mathrm{d}\bm{X}),
  \qquad
  \nabla_{\bm{X}}\log\det\bm{X}
  =
  \bm{X}^{-\top}.
\label{eq:analysis-logdet-gradient}
\end{equation}
正定矩阵对称，所以$\bm{X}^{-\top}=\bm{X}^{-1}$；对一般可逆实矩阵， $\log\det\bm{X}$还需要$\det\bm{X}>0$，而梯度应保留转置。以上推导依赖的不是元素级记忆口诀，而是局部线性近似、乘积次序和迹内积三件事。
\section{二阶导数与曲率}
\label{sec:analysis-second-order-curvature}
\subsection{Hessian矩阵描述梯度怎样变化}

一阶近似只保留函数的局部斜率。若还要判断这个斜率怎样随位置改变，就需要
\term{Hessian矩阵}（Hessian matrix）。设$f:\mathbb{R}^d\to\mathbb{R}$在 $\bm{x}$附近二阶可微，定义
\[   \nabla^2 f(\bm{x})   =
\begin{bmatrix}     \dfrac{\partial^2 f}{\partial x_1^2} &     \cdots &     \dfrac{\partial^2 f}{\partial x_1\partial x_d}\\     \vdots & \ddots & \vdots\\     \dfrac{\partial^2 f}{\partial x_d\partial x_1} &     \cdots &     \dfrac{\partial^2 f}{\partial x_d^2}
\end{bmatrix}. \]
若二阶偏导数在邻域内连续，混合偏导数可以交换，Hessian矩阵因而对称。对方向 $\bm{v}$，二次型 $\bm{v}^{\top}\nabla^2f(\bm{x})\bm{v}$是沿直线 $t\mapsto f(\bm{x}+t\bm{v})$在$t=0$处的二阶导数。它衡量的不是函数值大小，而是沿该方向的一阶斜率变化。

对二次函数
\[   f(\bm{x})   =   \frac12\bm{x}^{\top}\bm{A}\bm{x}   +\bm{b}^{\top}\bm{x}+c, \]
只有$\bm{A}$的对称部分影响函数值。令 $\bm{H}=(\bm{A}+\bm{A}^{\top})/2$，则 $\nabla f(\bm{x})=\bm{H}\bm{x}+\bm{b}$且 $\nabla^2f(\bm{x})=\bm{H}$。正特征值对应向上弯曲的方向，负特征值对应向下弯曲的方向；特征值绝对值相差很大时，不同方向的尺度也相差很大。
\subsection{Taylor展开把局部曲率写入余项}

设连接$\bm{x}$与$\bm{x}+\bm{h}$的线段包含在一个开集内，并且$f$在该开集上二阶连续可微。对一元函数$\phi(t)=f(\bm{x}+t\bm{h})$使用Taylor定理，存在 $\theta\in(0,1)$使
\begin{equation}
  f(\bm{x}+\bm{h})
  =
  f(\bm{x})
  +\nabla f(\bm{x})^{\top}\bm{h}
  +\frac12
  \bm{h}^{\top}
  \nabla^2f(\bm{x}+\theta\bm{h})
  \bm{h}.
\label{eq:analysis-second-order-taylor}
\end{equation}
若只假设Hessian矩阵沿线段连续，也可以使用不挑选中间点的积分余项
\[   f(\bm{x}+\bm{h})   =   f(\bm{x})   +\nabla f(\bm{x})^{\top}\bm{h}   +   \int_0^1(1-t)   \bm{h}^{\top}\nabla^2f(\bm{x}+t\bm{h})\bm{h}\,\mathrm{d}t. \]
于是当$\bm{h}\to\bm{0}$时，
\begin{equation}
  f(\bm{x}+\bm{h})
  =
  f(\bm{x})
  +\nabla f(\bm{x})^{\top}\bm{h}
  +\frac12\bm{h}^{\top}\nabla^2f(\bm{x})\bm{h}
  +o(\lVert\bm{h}\rVert_2^2).
\label{eq:analysis-second-order-local-model}
\end{equation}
式~\eqref{eq:analysis-second-order-local-model}是局部结论。若要在一个区域内给出统一余项界，还需控制该区域上的Hessian矩阵连续性或Lipschitz常数，不能仅凭一点处的二阶导数。
\subsection{二阶判据区分极值与鞍点}

设$\bm{x}^{\star}$是开集内部的局部极小点，且$f$在该点可微。对每个方向 $\bm{v}$考察$f(\bm{x}^{\star}+t\bm{v})$，一元函数的必要条件给出 $\nabla f(\bm{x}^{\star})^{\top}\bm{v}=0$。这对所有$\bm{v}$成立，因此
\begin{equation}
  \nabla f(\bm{x}^{\star})=\bm{0}.
\label{eq:analysis-first-order-necessary}
\end{equation}
若$f$还在该点二阶可微，同样沿任意方向使用一元二阶必要条件，可得
\begin{equation}
  \bm{v}^{\top}\nabla^2f(\bm{x}^{\star})\bm{v}\geq0
  \quad\text{对所有 }\bm{v},
\label{eq:analysis-second-order-necessary}
\end{equation}
即Hessian矩阵半正定。局部极大点对应半负定。这些只是必要条件：
$f(x)=x^4$与$f(x)=-x^4$在$0$处分别有严格局部极小值和严格局部极大值，
二者的二阶导数却都为零；$f(x)=x^3$在$0$处的一阶导数和二阶导数也都为零，
却没有极值。Hessian矩阵半正定或半负定本身不能决定驻点类型。

正定Hessian矩阵给出更强的充分条件。若$f$在$\bm{x}^{\star}$的某个邻域内二阶连续
可微，$\nabla f(\bm{x}^{\star})=\bm{0}$且
$\nabla^2f(\bm{x}^{\star})$正定，则其最小特征值为某个$\lambda>0$。由Hessian矩阵的
连续性，充分靠近$\bm{x}^{\star}$时，任意方向上的二次型仍至少为
$(\lambda/2)\lVert\bm{h}\rVert_2^2$。代入Taylor积分余项可得
\[   f(\bm{x}^{\star}+\bm{h})-f(\bm{x}^{\star})   \geq   \frac{\lambda}{4}\lVert\bm{h}\rVert_2^2>0 \]
对充分小的非零$\bm{h}$成立，故$\bm{x}^{\star}$是严格局部极小点。负定情形给出严格局部极大点。

若驻点处Hessian矩阵既有正特征值又有负特征值，沿相应特征向量的小扰动会使函数值分别增加和减小，该点称为\term{鞍点}（saddle point）。例如 $f(x,y)=x^2-y^2$的原点就是鞍点。若Hessian矩阵仅半正定或含零特征值，二阶判据可能没有结论，必须查看更高阶项或函数的其他结构。上述必要条件还假定点位于无约束开集内部； 边界极小点的梯度可以不为零，受约束问题需要后续章节的可行方向与乘子条件。
\section{积分与平均}
\label{sec:analysis-integration-average}
\subsection{测度先规定哪些集合能够赋予大小}

有限加权和把每个位置的函数值乘以一个权重再相加。积分把“位置”推广成可能不可数的集合，因此要先说明哪些集合可被测量，以及它们各有多大。集合$X$上的
\term{$\sigma$代数}（sigma-algebra）$\mathcal{A}$是一族子集，满足
\[   X\in\mathcal{A},\qquad   A\in\mathcal{A}\Longrightarrow X\setminus A\in\mathcal{A},\qquad   A_1,A_2,\ldots\in\mathcal{A}   \Longrightarrow   \bigcup_{n=1}^{\infty}A_n\in\mathcal{A}. \]
这些条件还蕴含对可数交封闭。$\mathcal{A}$中的集合称为可测集合。只要求可数次运算而非任意次运算，是因为极限与可数序列相连，同时又能避免给所有过于复杂的子集强行赋予相容的长度。

\term{测度}（measure）是映射$\mu:\mathcal{A}\to[0,\infty]$，满足 $\mu(\varnothing)=0$，并且对两两不交的$A_1,A_2,\ldots$有
\[   \mu\left(\bigcup_{n=1}^{\infty}A_n\right)   =   \sum_{n=1}^{\infty}\mu(A_n). \]
三元组$(X,\mathcal{A},\mu)$称为\term{测度空间}（measure space）。实数上的 Lebesgue测度把区间$(a,b)$赋值为$b-a$；有限集合上的计数测度把集合赋值为元素个数； 给有限点$x_i$分配质量$w_i\geq0$则得到 $\mu(A)=\sum_{i:x_i\in A}w_i$。因此有限加权和、通常的面积积分和后续的概率期望可以使用同一套记号
\scite{math-folland1999real}

在拓扑空间中，由所有开集生成的最小$\sigma$代数称为\term{Borel $\sigma$代数}
（Borel sigma-algebra）。实数上的开区间因而都是Borel可测集。Lebesgue测度可先由
区间覆盖定义外部大小，再用可测性判据扩张长度；把所有Lebesgue零测集的子集一并加入
Borel $\sigma$代数，得到相对于该测度的完备化。这里“完备”表示零测集的任意子集仍可测，
不同于第~\ref{sec:analysis-complete-contraction}节中“柯西数列不会逃出空间”的完备性。

可数可加性还给出集合极限与测度之间的连续性。若
$A_1\subseteq A_2\subseteq\cdots$且$A=\bigcup_nA_n$，令
$B_1=A_1$、$B_n=A_n\setminus A_{n-1}$，则$B_n$两两不交且
$A_n=\bigcup_{k=1}^nB_k$。因此
\[
  \mu(A)
  =
  \sum_{k=1}^{\infty}\mu(B_k)
  =
  \lim_{n\to\infty}\mu(A_n).
\]
若$A_1\supseteq A_2\supseteq\cdots$且$\mu(A_1)<\infty$，对递增集合
$A_1\setminus A_n$应用上一结论可得
$\mu(\bigcap_nA_n)=\lim_n\mu(A_n)$。递减情形需要$\mu(A_1)<\infty$：
若每个$A_n$的测度都为无穷，就不能通过“无穷减无穷”得到结论。例如Lebesgue测度下
$A_n=(n,\infty)$递减到空集，但每个$A_n$的测度仍为无穷。

若$\mu(N)=0$，则$N$称为\term{零测集}（null set）。某性质除一个零测集外处处成立， 就称它\term{几乎处处}（almost everywhere）成立，常缩写为a.e.。这不是说例外集合为空，而是说当前测度不给它任何质量；换用另一个测度后，原来的例外可能不再可忽略。
\subsection{可测函数保证阈值事件可以测量}

函数$f:X\to\overline{\mathbb{R}}=[-\infty,\infty]$称为\term{可测函数}
（measurable function），若每个$a\in\mathbb{R}$对应的集合
\[   \{x\in X:f(x)>a\} \]
都属于$\mathcal{A}$。这保证按函数值设置阈值产生的集合可以赋予测度。可测函数列的
逐点极限，以及可数族可测函数的逐点上确界和下确界，仍是可测函数。连续实函数相对于
Borel $\sigma$代数可测，但可测性比连续性弱得多。

形如
\[   s(x)=\sum_{j=1}^{m}a_j\mathbb{I}_{A_j}(x),   \qquad a_j\geq0,\quad A_j\in\mathcal{A}, \]
且$A_j$两两不交的函数称为非负\term{简单函数}（simple function）。它只取有限多个值，其积分自然定义为
\begin{equation}
  \int_X s\,\mathrm{d}\mu
  =
  \sum_{j=1}^{m}a_j\mu(A_j).
\label{eq:analysis-simple-integral}
\end{equation}
即使某个$\mu(A_j)=\infty$，该定义仍允许积分取$+\infty$，并约定 $0\cdot\infty=0$。
同一个简单函数可能有不同的分组表示；把这些表示按所有取值集合的交集细分成共同的
两两不交集合后，有限可加性给出相同的和，因此式~\eqref{eq:analysis-simple-integral}
不依赖具体表示。

每个非负可测函数$f$都能由一列非负简单函数从下方逐点逼近。对$n\geq1$，一个明确的
构造是
\begin{equation}
  s_n(x)
  =
  \sum_{k=0}^{n\,2^n-1}
  \frac{k}{2^n}
  \mathbb{I}_{\{k\,2^{-n}\leq f(x)<(k+1)2^{-n}\}}
  +
  n\mathbb{I}_{\{f(x)\geq n\}}.
\label{eq:analysis-canonical-simple-approximation}
\end{equation}
它先把函数值截断在$n$，再向下舍入到间距为$2^{-n}$的格点。截断高度随$n$增加，
格点也逐次细化，所以$0\leq s_n\leq s_{n+1}\leq f$。若$f(x)<\infty$，当
$n>f(x)$时还有$0\leq f(x)-s_n(x)<2^{-n}$；若$f(x)=+\infty$，则
$s_n(x)=n$。因此$s_n(x)\uparrow f(x)$对每个$x$成立。非负可测函数的积分定义为
\begin{equation}
  \int_X f\,\mathrm{d}\mu
  =
  \sup\left\{
    \int_X s\,\mathrm{d}\mu:
    0\leq s\leq f,\ s\text{为简单函数}
  \right\}.
\label{eq:analysis-nonnegative-lebesgue-integral}
\end{equation}
这里允许结果为$+\infty$。这个定义不要求把定义域切成连续小区间，而是按函数值分层， 所以能自然处理极限和零测例外。

\begin{lemma}[非负积分的线性与单调性]
\label{lem:analysis-nonnegative-integral-properties}
设$f,g:X\to[0,\infty]$可测，$\alpha,\beta\geq0$。若$f\leq g$几乎处处，则
\[
  \int_X f\,\mathrm{d}\mu\leq\int_X g\,\mathrm{d}\mu.
\]
此外，在$0\cdot\infty=0$的约定下，
\[
  \int_X(\alpha f+\beta g)\,\mathrm{d}\mu
  =
  \alpha\int_X f\,\mathrm{d}\mu
  +
  \beta\int_X g\,\mathrm{d}\mu.
\]
\end{lemma}
\begin{proof}
对非负简单函数，把两种表示细分到同一组两两不交的可测集合上，式
\eqref{eq:analysis-simple-integral}就退化为有限个非负数的求和，因而具有单调性和
非负线性。简单函数在可测零测集上的部分积分为零，所以由积分定义可知，两个可测非负
函数若只在一个可测零测集上取值不同，其积分相同。这里修改后的函数仍须可测；在未完备
测度空间中，不能对零测集作任意修改而省略这项要求。

对$f\leq g$几乎处处的情形，集合$N=\{x:f(x)>g(x)\}$可测且$\mu(N)=0$。把$f$在
$N$上改为$0$得到可测函数$\widetilde f$，其积分与$f$相同，并且逐点满足
$\widetilde f\leq g$。此时每个满足$0\leq s\leq\widetilde f$的简单函数也满足
$s\leq g$；
在式~\eqref{eq:analysis-nonnegative-lebesgue-integral}中取上确界便得到单调性。

还需把简单函数的线性提升到一般非负函数。先注意，若非负简单函数
$u_n\uparrow u$，则$\int u_n\,\mathrm{d}\mu\uparrow\int u\,\mathrm{d}\mu$。
单调性已经给出“$\leq$”方向。反过来，任取
$s=\sum_{j=1}^m a_j\mathbb{I}_{A_j}\leq u$，其中$a_j>0$且$A_j$两两不交，并固定
$c\in(0,1)$。集合
\[
  A_{j,n}=A_j\cap\{x:u_n(x)\geq ca_j\}
\]
随$n$递增到$A_j$，所以测度的下连续性给出
\[
  \lim_{n\to\infty}\int_Xu_n\,\mathrm{d}\mu
  \geq
  c\sum_{j=1}^m a_j\mu(A_j)
  =
  c\int_Xs\,\mathrm{d}\mu.
\]
令$c\uparrow1$，再对所有$s\leq u$取上确界，就得到所需等式。现在分别取
非负简单函数$s_n\uparrow f$和$t_n\uparrow g$，则
$\alpha s_n+\beta t_n\uparrow\alpha f+\beta g$。应用刚证明的简单函数单调极限性质
和简单函数的线性性，即得
\[
  \int_X(\alpha f+\beta g)\,\mathrm{d}\mu
  =
  \lim_{n\to\infty}
  \left(
    \alpha\int_Xs_n\,\mathrm{d}\mu
    +
    \beta\int_Xt_n\,\mathrm{d}\mu
  \right),
\]
右侧正是结论中的两个积分之和。
\end{proof}
\subsection{有符号积分需要控制正负两部分}

对可测实函数$f$，定义
\[   f^+(x)=\max\{f(x),0\},\qquad   f^-(x)=\max\{-f(x),0\}. \]
于是$f=f^+-f^-$且$|f|=f^++f^-$。只要 $\int f^+\,\mathrm{d}\mu$与$\int f^-\,\mathrm{d}\mu$不同时为无穷，就可以用两者之差定义$\int f\,\mathrm{d}\mu$。若
\begin{equation}
  \int_X|f|\,\mathrm{d}\mu<\infty,
\label{eq:analysis-absolute-integrability}
\end{equation}
就称$f$\term{绝对可积}（absolutely integrable）或Lebesgue可积。此时正负部分都有限，积分没有“无穷减无穷”的歧义，并满足
\[   \left|\int_X f\,\mathrm{d}\mu\right|   \leq   \int_X|f|\,\mathrm{d}\mu. \]

在闭区间上有界且Riemann可积的函数也Lebesgue可积，两种积分取值相同。Lebesgue积分
并非否定熟悉的面积解释，而是扩大可处理的函数，并让两个可测且几乎处处相等的函数
具有相同积分。例如在Lebesgue测度下，修改可测函数在单点上的值仍得到可测函数，也不
改变积分。反过来，积分有限不表示函数处处有界，也不自动给出连续性。
\subsection{乘积测度描述两个坐标上的平均}

给定两个测度空间$(X,\mathcal{A},\mu)$与$(Y,\mathcal{B},\nu)$，乘积 $\sigma$代数$\mathcal{A}\otimes\mathcal{B}$由矩形集合$A\times B$生成。在常用的$\sigma$有限条件下，存在唯一的\term{乘积测度}（product measure） $\mu\otimes\nu$满足
\[   (\mu\otimes\nu)(A\times B)=\mu(A)\nu(B). \]
对矩形上的指示函数，先对$x$积分或先对$y$积分都得到同一个乘积。一般函数能否继续交换积分次序，则取决于非负性或绝对可积性，第~\ref{sec:analysis-limit-integral-interchange}
节会给出准确条件。这里的$\sigma$有限是指空间能写成可数个有限测度集合之并；若脱离这一常见框架，乘积测度的唯一性和相应积分结论需要更谨慎的版本。

有限网格把这个构造还原为双重求和。若$X=\{x_1,\ldots,x_m\}$、
$Y=\{y_1,\ldots,y_n\}$分别带权重$u_i,v_j$，则
\[
  \int_{X\times Y}f\,\mathrm{d}(\mu\otimes\nu)
  =
  \sum_{i=1}^m\sum_{j=1}^n f(x_i,y_j)u_iv_j.
\]
有限和可以直接交换次序；一般测度空间中的累次积分还要处理无穷值与极限，
不能只由矩形公式推出。
\section{坐标变换及其积分应用}
\label{sec:analysis-coordinate-change-integral-applications}
\subsection{Jacobian矩阵的行列式衡量局部体积变化}

一维换元中，$u=T(x)$把长度$\mathrm{d}x$缩放为约 $|T'(x)|\mathrm{d}x$。在$d$维，局部线性近似
\[   T(\bm{x}+\bm{h})   \approx   T(\bm{x})+\bm{J}_T(\bm{x})\bm{h} \]
把一个小平行多面体的有向体积乘以 $\det\bm{J}_T(\bm{x})$，普通体积则乘以其绝对值。行列式为零意味着至少一个方向被压扁，局部逆映射通常不再存在。

设$U,V\subseteq\mathbb{R}^d$为开集， $T:U\to V$是连续可微双射，逆映射也连续可微，并且 $\det\bm{J}_T(\bm{x})\neq0$。对非负可测函数或绝对可积函数$g$，多维换元公式为
\begin{equation}
  \int_V g(\bm{y})\,\mathrm{d}\bm{y}
  =
  \int_U
  g\bigl(T(\bm{x})\bigr)
  \left|\det\bm{J}_T(\bm{x})\right|
  \,\mathrm{d}\bm{x}.
\label{eq:analysis-multivariate-change-variables}
\end{equation}
绝对值不可省略，因为积分中的体积不随坐标方向翻转而变成负数。一维单调换元是 $d=1$的特例；若在定积分中保留有向上下限，也可让导数符号由上下限次序吸收。
公式的证明可理解为把$U$分成足够小的区域：在每一块上用
$T(\bm{x}+\bm{h})\approx T(\bm{x})+\bm{J}_T(\bm{x})\bm{h}$替代原映射，
线性映射把该块体积乘以$|\det\bm{J}_T(\bm{x})|$，再让分割尺度趋于零。
这条路线只提供直观框架；严格证明还要控制局部近似误差、像集重叠和可测逼近，
不能把逐点线性化直接当作全局换元等式。
\begin{example}[极坐标中的面积权重]
\label{ex:analysis-polar-coordinate-jacobian}
令$D=\{\bm{x}\in\mathbb{R}^2:\lVert\bm{x}\rVert_2\leq1\}$，并令
$T(r,\theta)=(r\cos\theta,r\sin\theta)$。在开集
$(0,1)\times(0,2\pi)$上，$T$是到开圆盘内部去掉原点和正$x$轴径向接缝后的区域的
连续可微双射，逆映射也连续可微。其Jacobian矩阵的行列式为
\[   \det
\begin{bmatrix}     \cos\theta & -r\sin\theta\\     \sin\theta & r\cos\theta
\end{bmatrix}   =r. \]
圆周、原点和径向接缝都是二维零测集；参数矩形的边界也是零测集。因此在上述开集上应用
换元公式，再补回这些零测部分，便得到
\[   \int_D1\,\mathrm{d}\bm{x}   =   \int_0^{2\pi}\int_0^1r\,\mathrm{d}r\,\mathrm{d}\theta   =\pi. \]
这里积分上下限包含端点不改变积分值。若参数化在正测度区域上重复覆盖，则不能直接使用
上述双射版本，否则同一点会被重复计数。
\end{example}
\subsection{分部积分把导数从一个因子移到另一个因子}

若$u,v\in C^1([a,b])$，由乘积求导$(uv)'=u'v+uv'$并在$[a,b]$上积分，可得
\begin{equation}
  \int_a^b u(x)v'(x)\,\mathrm{d}x
  =
  [u(x)v(x)]_a^b
  -
  \int_a^b u'(x)v(x)\,\mathrm{d}x.
\label{eq:analysis-integration-by-parts}
\end{equation}
例如，取$u(x)=x$、$v'(x)=e^x$，得到 $\int_0^1xe^x\,\mathrm{d}x=[xe^x]_0^1-\int_0^1e^x\,\mathrm{d}x=1$。边界项不能无故删除；只有在函数于边界消失、区域无边界，或衰减速度足以使无穷远边界项趋零时才能省略。更一般地，若$u,v$绝对连续，公式仍成立，此时导数按几乎处处意义理解。
\subsection{路径积分和散度定理连接局部变化与整体边界}

设分段连续可微路径$\gamma:[a,b]\to\mathbb{R}^d$。对沿路径连续的标量函数
$h:\mathbb{R}^d\to\mathbb{R}$，标量线积分定义为
\[
  \int_\gamma h\,\mathrm{d}s
  =
  \int_a^b h(\gamma(t))\lVert\gamma'(t)\rVert_2\,\mathrm{d}t.
\]
速度的范数把参数增量换成弧长，因此反转路径方向不改变该积分。对沿路径连续的向量场
$\bm{F}:\mathbb{R}^d\to\mathbb{R}^d$，向量场线积分则定义为
\[   \int_{\gamma}\bm{F}\mathbin{\cdot}\mathrm{d}\bm{r}   =   \int_a^b   \bm{F}(\gamma(t))^{\top}\gamma'(t)\,\mathrm{d}t. \]
若$\bm{F}=\nabla\phi$且$\phi$连续可微，链式法则给出
\begin{equation}
  \int_{\gamma}\nabla\phi\mathbin{\cdot}\mathrm{d}\bm{r}
  =
  \phi(\gamma(b))-\phi(\gamma(a)).
\label{eq:analysis-line-integral-gradient}
\end{equation}
因此梯度场的路径积分只依赖端点。向量场线积分在路径方向反转时变号；若漏掉参数化导数
$\gamma'(t)$，就会丢失方向和行进速度的信息。

对具有分段光滑边界的有界区域$\Omega\subset\mathbb{R}^d$，若向量场$\bm{F}$
在$\overline{\Omega}$的一个开邻域内连续可微，
\term{散度定理}（divergence theorem）写成
\begin{equation}
  \int_{\Omega}\nabla\mathbin{\cdot}\bm{F}\,\mathrm{d}\bm{x}
  =
  \int_{\partial\Omega}\bm{F}^{\top}\bm{n}\,\mathrm{d}S,
\label{eq:analysis-divergence-theorem}
\end{equation}
其中$\bm{n}$是外法向量。左边汇总区域内部的局部流出率，右边计算穿过边界的净通量。它的严格适用条件随区域与函数空间版本而异；带尖点、无界区域或弱导数时，需要使用相应推广并重新核对边界项。分部积分、路径积分与散度定理共同说明，局部导数进入整体积分时， 通常会留下不可忽略的边界信息。
\section{极限与积分的交换}
\label{sec:analysis-limit-integral-interchange}
\subsection{逐点收敛不足以控制积分}

在$(0,1)$上令
\[   f_n(x)=n\mathbb{I}_{(0,1/n)}(x). \]
对每个固定$x>0$，充分大的$n$满足$x\geq1/n$，故$f_n(x)\to0$；然而
\[   \int_0^1f_n(x)\,\mathrm{d}x=1 \]
对所有$n$成立。因此 $\lim_n\int f_n\neq\int\lim_n f_n$。峰越来越窄，却同时越来越高，逐点观察不到被压缩区域中保留下来的总质量。要交换极限与积分，必须阻止这种质量逃逸或集中
\scite{math-folland1999real}
\subsection{单调收敛允许非负质量逐步增加}
\begin{theorem}[单调收敛定理]
\label{thm:analysis-monotone-convergence}
设$f_n,f:X\to[0,\infty]$可测，且$f_n(x)\uparrow f(x)$几乎处处，则
\[   \lim_{n\to\infty}\int_X f_n\,\mathrm{d}\mu   =   \int_X f\,\mathrm{d}\mu. \]
两边都允许取$+\infty$。
\end{theorem}
\begin{proof}
把单调性或收敛关系失效的可测零测集并为$N$，并将$f_n$与$f$在$N$上都改为$0$。
修改后的函数仍可测且积分不变，所以可假定$f_n\uparrow f$逐点成立。由
$0\leq f_n\leq f$，积分序列单调不减且
\[   L:=\lim_{n\to\infty}\int f_n\,\mathrm{d}\mu   \leq\int f\,\mathrm{d}\mu. \]
还需证明反向不等式。任取满足$0\leq s\leq f$的简单函数，将它写为 $s=\sum_{j=1}^m a_j\mathbb{I}_{E_j}$，其中$a_j>0$且$E_j$两两不交。固定 $c\in(0,1)$，定义
\[   E_{j,n}   =   E_j\cap\{x:f_n(x)\geq ca_j\}. \]
由于$f_n\uparrow f$且$f\geq a_j$于$E_j$，有$E_{j,n}\uparrow E_j$。测度对递增集合连续，故$\mu(E_{j,n})\uparrow\mu(E_j)$。同时
\[   f_n   \geq   c\sum_{j=1}^m a_j\mathbb{I}_{E_{j,n}}, \]
从而
\[   L   \geq   \lim_{n\to\infty}   c\sum_{j=1}^m a_j\mu(E_{j,n})   =   c\int s\,\mathrm{d}\mu. \]
先令$c\uparrow1$，再对所有$0\leq s\leq f$的简单函数取上确界。根据非负积分的定义~\eqref{eq:analysis-nonnegative-lebesgue-integral}， $L\geq\int f\,\mathrm{d}\mu$。与前一不等式合并即得结论。
\end{proof}

非负性避免正负无穷相消，单调性则保证已经积累的质量不会在后续消失。定理不要求存在一个可积的统一上界，也不要求各$\int f_n$有限。
\subsection{Fatou引理保留下极限方向的不等式}
对扩展实数列$(a_n)$，尾部下确界$\inf_{k\geq n}a_k$随$n$单调不减，尾部上确界
$\sup_{k\geq n}a_k$随$n$单调不增。因此可在扩展实数中定义
\[
  \liminf_{n\to\infty}a_n
  =
  \lim_{n\to\infty}\inf_{k\geq n}a_k,
  \qquad
  \limsup_{n\to\infty}a_n
  =
  \lim_{n\to\infty}\sup_{k\geq n}a_k.
\]
函数列的下极限与上极限按每个$x$逐点定义。若$a_n$收敛，则二者都等于通常极限；
若不收敛，它们分别记录尾部最终能够达到的下界与上界。
\begin{lemma}[Fatou引理]
\label{lem:analysis-fatou}
若$f_n:X\to[0,\infty]$可测，则
\begin{equation}
  \int_X\liminf_{n\to\infty}f_n\,\mathrm{d}\mu
  \leq
  \liminf_{n\to\infty}\int_Xf_n\,\mathrm{d}\mu.
\label{eq:analysis-fatou}
\end{equation}
\end{lemma}
\begin{proof}
令$g_n=\inf_{k\geq n}f_k$。则$g_n$可测， $0\leq g_n\uparrow\liminf_k f_k$。由单调收敛定理，
\[   \int\liminf_k f_k\,\mathrm{d}\mu   =   \lim_{n\to\infty}\int g_n\,\mathrm{d}\mu. \]
对每个$k\geq n$都有$g_n\leq f_k$，故
\[   \int g_n\,\mathrm{d}\mu   \leq   \inf_{k\geq n}\int f_k\,\mathrm{d}\mu. \]
令$n\to\infty$便得到式~\eqref{eq:analysis-fatou}。
\end{proof}

Fatou引理通常只给一个方向，因为积分可能在极限中损失质量。对上述收缩峰$f_n$， 左边为$0$，右边为$1$，不等式严格成立。
\subsection{支配收敛用一个可积函数统一控制}
\begin{theorem}[支配收敛定理]
\label{thm:analysis-dominated-convergence}
设$f_n,f:X\to\mathbb{R}$可测，$f_n\to f$几乎处处，并且存在非负可积函数$g$使
$|f_n|\leq g$几乎处处对所有$n$成立，则$f$可积，而且
\[   \lim_{n\to\infty}\int_Xf_n\,\mathrm{d}\mu   =   \int_Xf\,\mathrm{d}\mu. \]
此外，$\int|f_n-f|\,\mathrm{d}\mu\to0$。
\end{theorem}
\begin{proof}
把收敛或控制条件失效的可测零测集的可数并记为$N$，并将$f_n,f,g$在$N$上都改为
$0$。这些函数仍可测且积分不变，因而下述关系可以按逐点意义使用。
由$|f_n|\leq g$并取逐点极限可得$|f|\leq g$几乎处处，所以$f$可积。函数$g+f_n$与$g-f_n$均非负。对它们分别使用Fatou引理，
\begin{align*}
  \int g+\int f
  &\leq
  \int g+\liminf_{n\to\infty}\int f_n,\\
  \int g-\int f
  &\leq
  \int g-\limsup_{n\to\infty}\int f_n.
\end{align*}
由于$\int g<\infty$，可以从两边减去它，得到
\[   \int f   \leq\liminf_n\int f_n   \leq\limsup_n\int f_n   \leq\int f. \]
故积分收敛到$\int f$。

最后，$|f_n-f|\to0$几乎处处，且 $|f_n-f|\leq2g$。再对非负函数$|f_n-f|$应用刚证明的积分交换结论，得到 $\int|f_n-f|\,\mathrm{d}\mu\to0$。
\end{proof}

收缩峰$f_n=n\mathbb{I}_{(0,1/n)}$不存在满足条件的可积控制函数。事实上，若 $g\geq f_n$对所有$n$成立，则对接近$0$的$x$可选$n$约为$1/x$，迫使$g(x)$至少具有$1/x$量级，而它在$0$附近不可积。这个例子精确指出失效的是统一可积控制，而不是逐点极限。
\subsection{Tonelli与Fubini控制积分次序}
\begin{theorem}[Tonelli与Fubini定理]
\label{thm:analysis-tonelli-fubini}
设$(X,\mathcal{A},\mu)$与$(Y,\mathcal{B},\nu)$为$\sigma$有限测度空间，且 $f:X\times Y\to\overline{\mathbb{R}}$对乘积$\sigma$代数可测。
\begin{itemize}

\item 若$f\geq0$，则每个截面$y\mapsto f(x,y)$与$x\mapsto f(x,y)$都可测，并且两个
  内层积分函数
  \[
    x\longmapsto\int_Yf(x,y)\,\mathrm{d}\nu(y),
    \qquad
    y\longmapsto\int_Xf(x,y)\,\mathrm{d}\mu(x)
  \]
  分别关于$x$与$y$可测，且
  \[
    \begin{aligned}
      \int_{X\times Y}f\,\mathrm{d}(\mu\otimes\nu)
      &=
      \int_X\left(\int_Yf(x,y)\,\mathrm{d}\nu(y)\right)\mathrm{d}\mu(x)\\
      &=
      \int_Y\left(\int_Xf(x,y)\,\mathrm{d}\mu(x)\right)\mathrm{d}\nu(y).
    \end{aligned}
  \]
  允许共同值为$+\infty$。这是Tonelli定理。
\item 若$\int_{X\times Y}|f|\,\mathrm{d}(\mu\otimes\nu)<\infty$，则
  $y\mapsto f(x,y)$对几乎每个$x$可积，$x\mapsto f(x,y)$对几乎每个$y$可积。
  把相应内层积分函数在例外集上置为$0$后，它们分别可积，上述等式对$f$成立并取有限值。
  这是Fubini定理。
\end{itemize}
\end{theorem}

证明路线与本节前面的构造一致。先对矩形指示函数
$\mathbb{I}_{A\times B}$核对等式，两边都等于$\mu(A)\nu(B)$，再由有限可加性推广到
有限个矩形生成的代数。随后在有限测度块上使用单调类定理，把结论推广到乘积
$\sigma$代数中的一般可测集合；$\sigma$有限性允许用递增的有限测度块覆盖整个乘积
空间。至此，结论对任意乘积可测集合的指示函数成立，再由线性性推广到非负简单函数。
任意非负可测$f$可取简单函数$s_n\uparrow f$，三处积分分别使用单调收敛定理，便得到
Tonelli定理。若$f$绝对可积，对$f^+$与$f^-$分别应用Tonelli；$|f|$的积分有限保证
两部分的累次积分几乎处处有限，二者相减得到Fubini定理。

这里非负性允许无穷值，绝对可积性则排除了依赖积分次序的条件相消。例如在
$\mathbb{N}\times\mathbb{N}$上使用计数测度时，Tonelli说明非负双重级数可以任意交换
求和次序；Fubini说明绝对收敛的双重级数也可以这样做。若既不非负也不绝对可积，
交换次序可能改变结果，甚至一个次序收敛而另一个发散。因此“把两个积分换一下”
不是排版操作，而是需要先核对函数符号或绝对积分的数学步骤。
\section{隐式求导与参数积分}
\label{sec:analysis-implicit-parameter-integral}
\subsection{方程也能在局部定义函数}
显式关系$y=\varphi(x)$直接给出输入如何映射到输出，方程 $F(x,y)=0$却只规定二者必须共同满足的约束。例如 $F(x,y)=y^2-x=0$在$x>0$附近有$y=\sqrt{x}$和$y=-\sqrt{x}$两支，若不先指定所考察的解，方程并不定义唯一函数。更糟的是，在$(x,y)=(0,0)$附近，两支相交且 $\partial F/\partial y=2y=0$，不能把$y$稳定地看作$x$的可微函数。

下面的存在性证明需要用导数控制有限位移。对矩阵
$\bm{A}\in\mathbb{R}^{s\times r}$，由欧氏范数\term{诱导的算子范数}
（induced operator norm）定义为
\[
  \lVert\bm{A}\rVert_{\mathrm{op}}
  =
  \sup_{\bm{h}\neq\bm{0}}
  \frac{\lVert\bm{A}\bm{h}\rVert_2}{\lVert\bm{h}\rVert_2}
  =
  \sup_{\lVert\bm{h}\rVert_2=1}\lVert\bm{A}\bm{h}\rVert_2.
\]
它表示线性映射对向量长度的最大放大倍数，定义立即给出
$\lVert\bm{A}\bm{h}\rVert_2
\leq\lVert\bm{A}\rVert_{\mathrm{op}}\lVert\bm{h}\rVert_2$。
若连续可微映射$G$的定义域包含$\bm{y}'$到$\bm{y}$的线段，中值积分公式进一步给出
\begin{equation}
  \begin{aligned}
  \lVert G(\bm{y})-G(\bm{y}')\rVert_2
  &\leq
  \sup_{0\leq t\leq1}
  \left\lVert
    \bm{D}G\bigl(\bm{y}'+t(\bm{y}-\bm{y}')\bigr)
  \right\rVert_{\mathrm{op}}
  \lVert\bm{y}-\bm{y}'\rVert_2.
  \end{aligned}
\label{eq:analysis-operator-mean-value-bound}
\end{equation}
因此，邻域内统一的导数范数上界会直接变成Lipschitz常数；若该上界小于$1$，就得到
后文所需的压缩条件。
\begin{theorem}[隐函数定理]
\label{thm:analysis-implicit-function}
设$F:\mathbb{R}^p\times\mathbb{R}^m\to\mathbb{R}^m$在 $(\bm{a},\bm{b})$的某个开邻域内连续可微， $F(\bm{a},\bm{b})=\bm{0}$，并且关于待解变量的Jacobian矩阵 \[   \bm{D}_{\bm{y}}F(\bm{a},\bm{b})   \in\mathbb{R}^{m\times m} \]
可逆。则存在$\bm{a}$与$\bm{b}$的邻域$U,V$，以及唯一的连续可微函数
$\varphi:U\to V$，使$\varphi(\bm{a})=\bm{b}$，并且对所有
$(\bm{x},\bm{y})\in U\times V$都有
\[
  F(\bm{x},\bm{y})=\bm{0}
  \quad\Longleftrightarrow\quad
  \bm{y}=\varphi(\bm{x}).
\]
其Jacobian矩阵满足
\begin{equation}
  \bm{D}\varphi(\bm{x})
  =
  -
  \bigl[\bm{D}_{\bm{y}}F(\bm{x},\varphi(\bm{x}))\bigr]^{-1}
  \bm{D}_{\bm{x}}F(\bm{x},\varphi(\bm{x})).
\label{eq:analysis-implicit-derivative}
\end{equation}
\end{theorem}

存在性、局部唯一性和可微性都需要证明，不能由形式求导代替。第~\ref{sec:analysis-complete-contraction}节建立压缩映射后将补全证明。先假定定理给出的$\varphi$已经存在，对恒等式 $F(\bm{x},\varphi(\bm{x}))=\bm{0}$应用链式法则，得到
\[   \bm{D}_{\bm{x}}F   +   \bm{D}_{\bm{y}}F\,\bm{D}\varphi   =   \bm{0}. \]
左乘$\bm{D}_{\bm{y}}F$的逆即得式~\eqref{eq:analysis-implicit-derivative}。乘法次序由矩阵维度决定，不能把逆矩阵随意移到右侧。
\subsection{驻点方程给出最优解的局部变化}

考虑依赖参数$\theta$的无约束目标$q(z,\theta)$。若局部最优解$z^\star(\theta)$位于定义域内部，它满足驻点方程
\[   F(\theta,z)   =   \frac{\partial q}{\partial z}(z,\theta)   =   0. \]
若$q$二阶连续可微，且 $\partial^2q/\partial z^2$在所考察解处非零，隐函数定理给出该驻点的局部唯一分支。这并不自动说明它是全局最优解；还要由曲率或其他结构核对最优性。
\begin{example}[参数二次目标中的解导数与包络导数]
\label{ex:analysis-parametric-quadratic}
对$\theta>-1$，令
\[   q(z,\theta)   =   \frac12(1+\theta)z^2-\theta z. \]
因为$\partial^2q/\partial z^2=1+\theta>0$，目标严格凸，唯一驻点就是全局最小点。驻点方程为
\[   F(\theta,z)=(1+\theta)z-\theta=0. \]
隐式求导给出
\[   \frac{\mathrm{d}z^\star}{\mathrm{d}\theta}   =   -\frac{\partial F/\partial\theta}{\partial F/\partial z}   =   -\frac{z^\star-1}{1+\theta}. \]
直接消元得到
\[   z^\star(\theta)=\frac{\theta}{1+\theta},   \qquad   \frac{\mathrm{d}z^\star}{\mathrm{d}\theta}   =   \frac{1}{(1+\theta)^2}, \]
两种计算一致。

最优值$V(\theta)=q(z^\star(\theta),\theta)$的链式求导为
\[   V'(\theta)   =   \frac{\partial q}{\partial z}   \frac{\mathrm{d}z^\star}{\mathrm{d}\theta}   +   \frac{\partial q}{\partial\theta}. \]
在内部驻点，第一项因$\partial q/\partial z=0$消失，所以
\begin{equation}
  V'(\theta)
  =
  \left.
  \frac{\partial q}{\partial\theta}
  \right|_{z=z^\star(\theta)}
  =
  \frac12(z^\star)^2-z^\star.
\label{eq:analysis-envelope-quadratic}
\end{equation}
另一方面，直接代回可得
\[   V(\theta)=-\frac{\theta^2}{2(1+\theta)},   \qquad   V'(\theta)   =   -\frac{\theta(\theta+2)}{2(1+\theta)^2}, \]
与式~\eqref{eq:analysis-envelope-quadratic}一致。
\end{example}

式~\eqref{eq:analysis-envelope-quadratic}体现了\term{包络定理}
（envelope theorem）的基本机制：最优值的一阶变化不需要显式乘上最优解导数，因为目标对决策变量的一阶导数在内部最优点为零。但解导数本身并未消失；若关心参数变化后 $z^\star$怎样移动，仍要使用隐式求导。边界解、不可微目标、多重最优解或受约束问题不能直接套用这个简单版本，后续最优化章节会加入相应条件。
\subsection{积分号下求导需要统一可积控制}

设
\[   G(\theta)=\int_X g(x,\theta)\,\mathrm{d}\mu(x). \]
形式上希望把导数移入积分号：
\[   G'(\theta)   \stackrel{?}{=}   \int_X\frac{\partial g}{\partial\theta}(x,\theta)\,\mathrm{d}\mu(x). \]
真正需要交换的是差商的极限与积分。固定$\theta_0$，设对每个$\theta\in I$，
$x\mapsto g(x,\theta)$可测，其中$I$是$\theta_0$的一个开邻域，并且
$g(\cdot,\theta_0)$可积。若对几乎处处的$x$，
$\theta\mapsto g(x,\theta)$在$I$内可微，偏导数关于$x$可测，并且存在可积函数$h$使
\begin{equation}
  \left|
    \frac{\partial g}{\partial\theta}(x,\theta)
  \right|
  \leq h(x)
  \quad
  \text{对几乎处处的$x$及所有$\theta\in I$成立},
\label{eq:analysis-parameter-derivative-domination}
\end{equation}
那么中值定理把差商的绝对值也控制在$h(x)$之下。差商逐点收敛到 $\partial g/\partial\theta(x,\theta_0)$，支配收敛定理于是给出
\begin{equation}
  G'(\theta_0)
  =
  \int_X
  \frac{\partial g}{\partial\theta}(x,\theta_0)
  \,\mathrm{d}\mu(x).
\label{eq:analysis-differentiate-under-integral}
\end{equation}
条件~\eqref{eq:analysis-parameter-derivative-domination}是一个便于核对的充分条件； 更一般的版本可以直接控制差商，或使用其他一致可积条件。
\begin{example}[参数积分中的控制函数失效]
\label{ex:analysis-parameter-integral-failure}
下面考察参数域边界处对应的右求导问题。对$\theta\geq0$与$x\in(0,1)$，定义
\[   g(x,\theta)   =
\begin{cases}     \dfrac{\theta x}{x^2+\theta^2}, & \theta>0,\\[5pt]     0, & \theta=0.
\end{cases} \]
每个固定$x>0$处关于$\theta$的右偏导数都等于$1/x$，但$1/x$在$(0,1)$上不可积。
直接积分得到
\[   G(\theta)   =   \frac{\theta}{2}   \log\left(\frac{1+\theta^2}{\theta^2}\right)   \quad(\theta>0),   \qquad G(0)=0. \]
于是$G(\theta)/\theta$在$\theta\downarrow0$时发散，$G$在$0$处没有有限右导数。积分内的
逐点右导数虽然存在，却缺少统一可积控制；因此这里不能交换右求导与积分。
\end{example}
\section{非光滑与人为替代梯度}
\label{sec:analysis-nonsmooth-surrogate-gradient}
\subsection{不可微点与几乎处处可导是不同陈述}

函数$|x|$在$x\neq0$处导数为$\operatorname{sign}(x)$，在$0$处左右导数不同。 ReLU函数$\max\{0,x\}$也只在$0$处不可微。由于单点是Lebesgue零测集，它们都几乎处处可导，但“几乎处处有导数”并没有为例外点指定导数值。

阈值函数$\mathbb{I}\{x\geq0\}$在$x\neq0$处导数为$0$，在$0$处不连续且不可微。若机械地使用几乎处处导数，反向传播得到的信号几乎处处为零，这无法表达跨过阈值会改变输出。重置、取整和离散选择也有类似问题：前向映射的真实局部导数可能不存在或不能提供有用方向。

凸函数的不可微点可以用凸次梯度处理，其正式定义与最优性条件见
第~\ref{sec:opt-nonsmooth-proximal}节。凸次梯度依赖全局仿射下界，并非“随便给不可微点填
一个导数”；对于一般非凸函数，还需要区分其他广义导数概念。

有限维\term{局部Lipschitz函数}是在每个点附近都满足某个Lipschitz界的函数。
Rademacher定理说明，这类函数几乎处处可微，但定理仍没有为例外点指定唯一梯度。
一种适用于局部Lipschitz函数的扩展是\term{Clarke广义梯度}
（Clarke generalized gradient）：
\[
  \partial^\circ f(\bm{x})
  =
  \overline{\operatorname{conv}}
  \left\{
    \lim_{k\to\infty}\nabla f(\bm{x}_k):
    \bm{x}_k\to\bm{x},\
    f\text{在}\bm{x}_k\text{处可微，且该极限存在}
  \right\}.
\]
这里$\operatorname{conv}$表示有限个非负权重和为$1$的凸组合，横线表示再取闭包。
对局部Lipschitz凸函数，Clarke广义梯度与
第~\ref{sec:opt-nonsmooth-proximal}节定义的凸次梯度集合一致；对一般非凸局部Lipschitz
函数，前者仍由邻近可微点的梯度极限构造，不再表示全局仿射下界。阈值函数在跳点附近
不连续，因而不是局部Lipschitz函数，不能直接套用这一构造。
\subsection{替代梯度是反向计算规则}

在含阈值或脉冲事件的模型中，常保留离散前向映射，却在反向计算时把其导数替换成某个平滑且非零的函数。这种信号称为\term{替代梯度}（surrogate gradient）。它是算法人为指定的近似更新方向，不自动等于前向映射的经典导数，也通常不满足凸次梯度的不等式。因此使用时应分别写清前向函数、反向替代函数及其尺度，不能把训练规则描述成对原离散函数的精确求导。

中心差分
\[   \frac{f(x+\varepsilon)-f(x-\varepsilon)}{2\varepsilon} \]
只能检验实际代入的函数$f$在尺度$\varepsilon$附近怎样变化。若程序反向传播使用替代梯度，而数值差分计算离散前向函数，两者本来就不是同一个导数对象；不一致不必然说明程序实现错误。
\subsection{函数值逼近不保证导数逼近}

用平滑函数
\[   f_{\varepsilon}(x)=\sqrt{x^2+\varepsilon^2},   \qquad \varepsilon>0 \]
逼近$|x|$。由
\[   0\leq f_{\varepsilon}(x)-|x|   =   \frac{\varepsilon^2}   {\sqrt{x^2+\varepsilon^2}+|x|}   \leq\varepsilon, \]
可知$f_{\varepsilon}$在整个$\mathbb{R}$上一致收敛到$|x|$。其导数为
\[   f_{\varepsilon}'(x)   =   \frac{x}{\sqrt{x^2+\varepsilon^2}}. \]
对每个$x\neq0$，导数收敛到$\operatorname{sign}(x)$；在$x=0$处始终为$0$，而$|x|$
在那里没有导数。即使人为规定$\operatorname{sign}(0)=0$，导数收敛也不一致。例如取
$x_\varepsilon=\varepsilon^2>0$，则
\[
  f_\varepsilon'(x_\varepsilon)
  =
  \frac{\varepsilon}{\sqrt{1+\varepsilon^2}}
  \longrightarrow0,
  \qquad
  \operatorname{sign}(x_\varepsilon)=1.
\]
因此两者之差趋于$1$，不可能一致收敛。所以函数值的一致逼近不能推出导数的一致逼近。
讨论光滑近似时，必须明确控制的是函数值、导数还是积分误差。
\section{函数收敛与本质界}
\label{sec:analysis-function-convergence-essential-bounds}
\subsection{同一函数列可以在不同误差尺度下得到不同结论}

设$f_n,f:X\to\mathbb{R}$。逐点收敛允许开始稳定的时刻依赖输入：
\[   \forall x,\quad |f_n(x)-f(x)|\to0. \]
一致收敛要求一个$n$同时控制所有输入：
\begin{equation}
  \lVert f_n-f\rVert_{\infty}
  :=
  \sup_{x\in X}|f_n(x)-f(x)|
  \longrightarrow0.
\label{eq:analysis-uniform-convergence}
\end{equation}
在测度空间上，$L^p$收敛则控制平均的$p$次误差：
\begin{equation}
  \lVert f_n-f\rVert_p
  :=
  \left(
    \int_X|f_n-f|^p\,\mathrm{d}\mu
  \right)^{1/p}
  \longrightarrow0,
  \qquad 1\leq p<\infty.
\label{eq:analysis-lp-convergence}
\end{equation}

在$[0,1]$上取$f_n(x)=x^n$。它逐点收敛到
\[   f(x)=
\begin{cases}     0, & 0\leq x<1,\\     1, & x=1.
\end{cases} \]
但收敛不一致，因为对每个$n$， $\sup_{0\leq x<1}x^n=1$。另一方面，对任意有限$p\geq1$，
\[   \lVert f_n-f\rVert_p^p   =   \int_0^1x^{np}\,\mathrm{d}x   =   \frac{1}{np+1}   \longrightarrow0. \]
同一函数列因此逐点且在每个有限$L^p$意义下收敛，却不一致收敛。区别来自最坏输入与平均误差采用了不同尺度。

若$\mu(X)<\infty$，一致收敛蕴含$L^p$收敛，因为
\[   \lVert f_n-f\rVert_p   \leq   \mu(X)^{1/p}\lVert f_n-f\rVert_{\infty}. \]
当空间测度无穷时，这个推论可能失败；常数函数$f_n=1/n$在$\mathbb{R}$上一致趋于 $0$，但它甚至不属于有限$p$的$L^p(\mathbb{R})$。
\subsection{本质上确界忽略零测例外}

普通上确界会被单个异常点改变，而保持可测性的零测集修改不影响积分。可测函数$f$的
\term{本质上确界}（essential supremum）定义为
\[   \operatorname*{ess\,sup}_{x\in X}|f(x)|   =   \inf\{M\geq0:|f(x)|\leq M\text{几乎处处}\}. \]
由此定义
\begin{equation}
  \lVert f\rVert_{L^\infty}
  =
  \operatorname*{ess\,sup}_{x\in X}|f(x)|.
\label{eq:analysis-essential-supremum-norm}
\end{equation}
例如在$[0,1]$上令$f(0)=100$、其余位置$f(x)=0$，则普通上确界为$100$， 本质上确界却为$0$。

对$1\leq p<\infty$，式~\eqref{eq:analysis-lp-convergence}中的表达式满足齐次性和三角不等式，但若只把逐点不同的函数当作不同对象，$\lVert f\rVert_p=0$仍不能推出 $f$处处为零。为得到真正的范数空间，定义
\[   f\sim g   \quad\Longleftrightarrow\quad   f=g\text{几乎处处}, \]
并把$L^p$定义为可测函数按此关系形成的等价类。$L^\infty$也采用同样约定。这与所有有界函数上的普通上确界范数不同：后者区分单点修改，前者不区分。
当$0<p<1$时，同一积分表达式仍可用来度量误差大小，但通常不满足三角不等式，
所以不是范数；本章不把这一区间记作范数空间。
\subsection{误差尺度必须与问题中的分布和函数类对应}

一个固定函数序列的一致收敛，讨论的是$n$增大时对所有输入能否共同控制。统计学习中“经验平均一致逼近总体平均”通常还要对候选函数类取上确界，控制对象形如
\[   \sup_{f\in\mathcal{F}}   \left|     \frac1n\sum_{i=1}^nf(X_i)-\int f\,\mathrm{d}P   \right|. \]
这里的一致性针对候选函数$f\in\mathcal{F}$，不是本节式~\eqref{eq:analysis-uniform-convergence}中针对输入$x\in X$的一致收敛。两者都出现上确界，却不能混为同一概念。

几乎处处和$L^p$范数还依赖所选测度。若$Q$对$P$\term{绝对连续}
（absolutely continuous），记作$Q\ll P$，则$P(A)=0$蕴含$Q(A)=0$，原来 $P$意义下的零测例外在$Q$下仍可忽略。若没有$Q\ll P$，换分布可能给原零测集正质量。例如$P$是$[0,1]$上的均匀分布，$Q$把全部质量放在$0$点；两个只在$0$处不同的函数在$P$意义下等价，在$Q$意义下却可能具有完全不同的误差。
\section{完备空间与压缩映射定理}
\label{sec:analysis-complete-contraction}
\subsection{完备性保证柯西极限留在空间内}

度量空间$(X,d)$若每个柯西数列都收敛到$X$中的点，就称为\term{完备}
（complete）。实数空间与有限维欧氏空间完备，有理数空间在通常距离下不完备。完备性不是说每个序列都收敛，而是说只要序列内部已经表现出应有的趋近性质，其极限不会落到空间外。

令$S$为任意非空集合，$\mathcal{B}(S)$为所有有界实函数$S\to\mathbb{R}$的集合， 配备普通上确界范数
\[   \lVert f\rVert_{\infty}   =   \sup_{x\in S}|f(x)|. \]
这个空间完备。事实上，若$(f_n)$在上确界范数下为柯西列，则对每个固定$x$， $(f_n(x))$是实柯西列，故存在$f(x)=\lim_n f_n(x)$。取一个$n_0$使 $\lVert f_n-f_{n_0}\rVert_\infty\leq1$对所有充分大的$n$成立，再令$n\to\infty$， 得到$|f(x)|\leq|f_{n_0}(x)|+1$，所以$f$有界。最后，对任意$\varepsilon>0$， 选$N$使$m,n\geq N$时 $\lVert f_m-f_n\rVert_\infty<\varepsilon$。固定$n\geq N$并令$m\to\infty$， 对所有$x$都有$|f(x)-f_n(x)|\leq\varepsilon$，于是 $\lVert f-f_n\rVert_\infty\leq\varepsilon$。这证明$f_n$在该空间中收敛到$f$。

空间$L^\infty$也完备，但其元素是几乎处处等价类，范数是式~\eqref{eq:analysis-essential-supremum-norm}。上面的逐点证明不能原封不动搬过去， 因为不同$n$的代表函数可能各自在不同零测集上修改；严格证明需要先统一这些例外集合。
\subsection{压缩映射把距离按固定比例缩短}

映射$T:X\to X$若存在$q\in[0,1)$使
\begin{equation}
  d(Tx,Ty)\leq qd(x,y)
  \quad\text{对所有 }x,y\in X,
\label{eq:analysis-contraction-definition}
\end{equation}
就称为\term{压缩映射}（contraction mapping）。关键是同一个$q<1$在整个所考察空间上成立，并且$T$确实把$X$映回自身
\scite{math-conway1990functional}
\begin{theorem}[Banach压缩映射定理]
\label{thm:analysis-banach-fixed-point}
设$(X,d)$是非空完备度量空间，$T:X\to X$为压缩映射，压缩因子为$q<1$。则$T$有
唯一不动点$x^\star\in X$。从任意$x^{(0)}\in X$出发定义
$x^{(n+1)}=T(x^{(n)})$，都有$x^{(n)}\to x^\star$，并且
\begin{align}
  d(x^{(n)},x^\star)
  &\leq
  \frac{q^n}{1-q}d(x^{(1)},x^{(0)}),
\label{eq:analysis-banach-apriori-error}\\
  d(x^{(n)},x^\star)
  &\leq
  \frac{q}{1-q}d(x^{(n)},x^{(n-1)})
  \quad(n\geq1).
\label{eq:analysis-banach-aposteriori-error}
\end{align}
\end{theorem}
\begin{proof}
反复使用压缩条件可得
\[   d(x^{(n+1)},x^{(n)})   \leq   q^n d(x^{(1)},x^{(0)}). \]
对$m>n$，三角不等式和几何级数给出
\begin{align*}
  d(x^{(m)},x^{(n)})
  &\leq
  \sum_{k=n}^{m-1}d(x^{(k+1)},x^{(k)})\\
  &\leq
  d(x^{(1)},x^{(0)})\sum_{k=n}^{m-1}q^k\\
  &\leq
  \frac{q^n}{1-q}d(x^{(1)},x^{(0)}).
\end{align*}
右边随$n\to\infty$趋于零，所以$(x^{(n)})$是柯西数列。完备性保证存在
$x^\star\in X$使$x^{(n)}\to x^\star$。压缩条件还蕴含$T$为Lipschitz连续，因此
\[   T(x^\star)   =   T\left(\lim_nx^{(n)}\right)   =   \lim_nT(x^{(n)})   =   \lim_nx^{(n+1)}   =   x^\star. \]
故$x^\star$是不动点。

若$y^\star$也是不动点，则
\[   d(x^\star,y^\star)   =   d(Tx^\star,Ty^\star)   \leq qd(x^\star,y^\star). \]
因为$q<1$，只能有$d(x^\star,y^\star)=0$，故不动点唯一。

在上面的柯西界中令$m\to\infty$，便得式~\eqref{eq:analysis-banach-apriori-error}。若从第$n-1$步重新计算同一界，
\[   d(x^{(n)},x^\star)   \leq   \sum_{k=n}^{\infty}d(x^{(k+1)},x^{(k)})   \leq   d(x^{(n)},x^{(n-1)})\sum_{j=1}^{\infty}q^j, \]
便得到式~\eqref{eq:analysis-banach-aposteriori-error}。
\end{proof}

这个证明同时说明存在、唯一和可计算的误差界各来自哪里：几何尾和使迭代为柯西列， 完备性把极限留在空间中，连续性使极限成为不动点，严格的$q<1$保证唯一性。若只允许$q=1$，结论都可能丢失。恒等映射$T(x)=x$有无穷多个不动点； $T(x)=x+1$在$\mathbb{R}$上没有不动点，且相应迭代发散。非扩张并不等于压缩。
完备性也不能省略。在非完备空间$X=(0,1)$上，映射$T(x)=x/2$是压缩因子为
$1/2$的自映射，但它唯一可能的不动点是被排除的$0$。从任意$x^{(0)}\in X$出发，
迭代$x^{(n)}=x^{(0)}/2^n$在$\mathbb{R}$中趋于$0$，却不在$X$中收敛。
\subsection{压缩迭代补全隐函数的局部存在证明}

现在回到定理~\ref{thm:analysis-implicit-function}。记
\[   \bm{A}   =   \bm{D}_{\bm{y}}F(\bm{a},\bm{b}),   \qquad   T_{\bm{x}}(\bm{y})   =   \bm{y}-\bm{A}^{-1}F(\bm{x},\bm{y}). \]
$F(\bm{x},\bm{y})=\bm{0}$当且仅当$\bm{y}$是$T_{\bm{x}}$的不动点。并且
\[   \bm{D}_{\bm{y}}T_{\bm{x}}(\bm{y})   =   \bm{I}   -   \bm{A}^{-1}\bm{D}_{\bm{y}}F(\bm{x},\bm{y}), \]
它在$(\bm{a},\bm{b})$处为零。由导数连续性，可以选择闭球 $B=\{\bm{y}:\lVert\bm{y}-\bm{b}\rVert_2\leq r\}$和$\bm{a}$的邻域$U$，使
\begin{equation}
  \sup_{\bm{x}\in U,\,\bm{y}\in B}
  \left\lVert
    \bm{D}_{\bm{y}}T_{\bm{x}}(\bm{y})
  \right\rVert_{\mathrm{op}}
  \leq q<1.
\label{eq:analysis-ift-uniform-contraction}
\end{equation}
沿$\bm{y}$方向使用式~\eqref{eq:analysis-operator-mean-value-bound}可知，每个
$T_{\bm{x}}$在$B$上都是同一因子$q$的压缩映射。

还必须核对自映射。因为$T_{\bm{a}}(\bm{b})=\bm{b}$且$T_{\bm{x}}(\bm{b})$关于 $\bm{x}$连续，可进一步缩小$U$，使
\[   \lVert T_{\bm{x}}(\bm{b})-\bm{b}\rVert_2   \leq\frac{1-q}{2}r. \]
于是对$\bm{y}\in B$，
\[
  \lVert T_{\bm{x}}(\bm{y})-\bm{b}\rVert_2
  \leq
  q\lVert\bm{y}-\bm{b}\rVert_2+\frac{1-q}{2}r
  \leq
  \frac{1+q}{2}r
  <r.
\]
所以$T_{\bm{x}}(B)$包含在$B$的内部$V$中。闭球$B$在有限维欧氏空间中完备，Banach
定理便对每个$\bm{x}\in U$给出唯一不动点$\varphi(\bm{x})\in V$。若
$(\bm{x},\bm{y})\in U\times V$且$F(\bm{x},\bm{y})=\bm{0}$，则$\bm{y}\in B$也是
$T_{\bm{x}}$的不动点，唯一性迫使$\bm{y}=\varphi(\bm{x})$。这证明了定理所述的局部
存在与零点集合的逐点唯一性。

还需说明所得解随参数连续可微。对$\bm{x},\bm{x}'\in U$，在更小的紧邻域内， $T_{\bm{x}}(\bm{y})$关于$\bm{x}$的导数有统一上界$L$。利用两个不动点关系，
\[   \lVert\varphi(\bm{x})-\varphi(\bm{x}')\rVert_2   \leq   q\lVert\varphi(\bm{x})-\varphi(\bm{x}')\rVert_2   +L\lVert\bm{x}-\bm{x}'\rVert_2, \]
故$\varphi$局部Lipschitz连续。令 $\bm{h}\to\bm{0}$并记 $\bm{k}=\varphi(\bm{x}+\bm{h})-\varphi(\bm{x})$，前述界给出 $\lVert\bm{k}\rVert_2=O(\lVert\bm{h}\rVert_2)$。对恒等式
\[   F(\bm{x}+\bm{h},\varphi(\bm{x})+\bm{k})   -   F(\bm{x},\varphi(\bm{x}))   =   \bm{0} \]
作一阶展开，得到
\[   \bm{D}_{\bm{x}}F\,\bm{h}   +   \bm{D}_{\bm{y}}F\,\bm{k}   +   o(\lVert\bm{h}\rVert_2)   =   \bm{0}. \]
邻域取得足够小时，$\bm{D}_{\bm{y}}F$仍可逆，于是
\[   \bm{k}   =   -   (\bm{D}_{\bm{y}}F)^{-1}   \bm{D}_{\bm{x}}F\,\bm{h}   +   o(\lVert\bm{h}\rVert_2). \]
这证明$\varphi$可微并给出式~\eqref{eq:analysis-implicit-derivative}；各Jacobian矩阵连续又保证$\bm{D}\varphi$连续。至此，隐函数定理中的存在、局部唯一与导数公式都已得到说明。证明中特别核对了自映射和统一压缩条件，只有形式上的固定点改写并不足够。
\section{本章小结}
\label{sec:analysis-summary}

无限过程的核心问题不是“项数足够多”或“变化足够小”，而是误差采用什么尺度、极限是否仍在所考察空间，以及交换运算需要哪些统一控制。数列的柯西条件和空间的完备性回答极限能否留在空间内；紧致性与连续性共同保证极值能够达到；可求和误差比单步误差趋零提供了更强的累计控制。

微分把局部变化组织成线性映射。梯度、Jacobian矩阵和矩阵梯度分别对应不同输入输出形状， 链式法则中的乘积次序由这些形状决定。Hessian矩阵与二阶Taylor展开描述局部曲率，但二阶判据有明确的光滑性、内点和正定性边界。隐函数定理进一步说明，可逆的待解变量Jacobian矩阵 如何把方程的局部解变成可微函数；最优解导数与最优值的包络导数由此得到区分。

积分把有限加权和推广到测度空间。简单函数从下逼近给出非负积分，正负部分和绝对可积性消除有符号积分的歧义。换元公式、分部积分与散度定理都表明局部微分进入整体积分时会带来体积权重或边界项。单调收敛、Fatou引理和支配收敛则准确规定极限何时能进入积分号； Tonelli依靠非负性，Fubini依靠绝对可积性来交换多重积分次序。

当函数不光滑时，几乎处处导数、凸次梯度、Clarke广义梯度与人为替代梯度是不同对象。函数值、导数、逐点误差、一致误差和$L^p$误差也不能互相替代；零测例外还依赖当前测度。最后，Banach 压缩映射定理把这些条件落实为可执行的迭代保证：统一压缩产生几何尾和，完备性给出存在， 严格压缩给出唯一性，并同时提供误差界。后续章节使用微分、积分或固定点时，应先核对本章列出的定义域、光滑性、可积性、完备性与统一控制条件，再进行形式计算。
第~\ref{chap:optimization}章将把这些局部变化、曲率与固定点工具用于最优性条件和求解算法，
并区分解的存在、唯一、稳定与可计算性。
